\chapter{Its a Bird! Its a PLANE!}\label{ch:aviation}
\markboth{Its a Bird! Its a Plane!}{Its a Bird! Its a Plane!}

\chapquote{``I can assure you that flying saucers, given that they exist, are not constructed by any power on Earth.''}{Harry S.~Truman}

% ---------------------------------------------------------------
\section{The History of Aviation}\label{sec:historyaviation}

A UFO is an \emph{unidentified flying object}, and identifying flying objects requires knowing
what is up there. Most people do not, and that single gap in general knowledge is responsible for
a large share of the reports in this book. So: a short history of what human beings have put in
the air, and when.

Balloons came first, and they matter more than anybody expects. The Montgolfier brothers flew a
hot-air balloon at Annonay in June 1783; Jacques Charles flew hydrogen in December of the same
year. Balloons were used for military observation at Fleurus in 1794 and extensively in the
American Civil War. By the late nineteenth century they were commonplace, which is worth
remembering when you read a nineteenth-century airship sighting.

\figwrap[r]{0.38}{Orville at the controls, Wilbur running alongside, Kill Devil Hills, 17 December 1903. Twelve seconds that the rest of this chapter depends on.}{https://commons.wikimedia.org/wiki/File:First_flight2.jpg}{fig:wright}
Heavier-than-air flight arrived on 17 December 1903 at Kill Devil Hills, North Carolina --- 12
seconds, 120 feet (see Figure~\ref{fig:wright}). What the Wrights actually contributed was not the wing but the control
problem: three-axis control, coupled roll and yaw, and their own wind tunnel to get usable lift
data after they found the published tables wrong. That is engineering of a high order and it is
why they succeeded where better-funded rivals failed.

Then the timeline compres, and the compression is the point. Powered flight to a monoplane
fighter is eleven years. To all-metal airliners, twenty-five. Frank Whittle and Hans von Ohain
developed turbojets independently in the 1930s; the Heinkel He 178 flew in 1939 and the Me 262
entered service in 1944. Chuck Yeager exceeded Mach 1 in the Bell X-1 in October 1947 --- the
same year as Roswell and Kenneth Arnold, which is not a coincidence and I will return to it.
Sputnik, 1957. Gagarin, 1961. Armstrong, 1969. Sixty-six years from a sand dune in North
Carolina to the Sea of Tranquility.

\figwrap[r]{0.38}{William Powell ``Bill'' Lear (1902--1978): car radios, autopilots, the 8-track cartridge, and the aircraft that put private owners above the airliners.}{https://commons.wikimedia.org/wiki/File:Bill_Lear.jpg}{fig:billlear}
One strand of that compression runs through a single family, and because it will matter twice more
in this book it is worth setting down here. William Powell Lear (1902--1978) was a self-taught
engineer with no formal training past the eighth grade and something over a hundred patents to his
name (see Figure~\ref{fig:billlear}). He built the first practical car radio and sold the business that became Motorola;
he built the lightweight aircraft radio direction finder and then the first dependable autopilot
small enough for a business aeroplane, which won him the Collier Trophy in 1950; and in 1965 he
commercialised the 8-track tape cartridge, which is why a generation associates his surname with
dashboards rather than with wings.

\figwrap[l]{0.38}{A Gates Learjet 23. First flight 7 October 1963: a private aircraft that cruised in the high forties, above the airliners and near the reconnaissance traffic.}{https://commons.wikimedia.org/wiki/File:Gates_Learjet_23_(10438289895).jpg}{fig:learjet23}
The aircraft is the part that belongs in a UFO book. The Learjet 23 first flew on 7 October 1963,
derived from a Swiss ground-attack design, and it certified in mid-1964 as the first mass-produced
business jet (see Figure~\ref{fig:learjet23}). It was small, very fast, and it cruised in the mid to high forties --- twenty
thousand feet above the piston airliners of the previous decade and in the same block of sky as
traffic the public had no idea existed. Bill Lear did not merely shorten the timeline above; he put
private citizens into the altitude band where the interesting sightings happen. His son would spend
thirty years flying in that band, and then spend the rest of his life telling the world what he
believed was hidden in it.

\figwrap[l]{0.38}{The SR-71 Blackbird: classified while flying, and unlike anything the public had a category for.}{https://commons.wikimedia.org/wiki/File:Lockheed_SR-71_Blackbird.jpg}{fig:sr71}
The SR-71 Blackbird first flew in 1964, cruised above Mach 3 at 85,\!000 feet, leaked fuel on the ground (see Figure~\ref{fig:sr71}) because its panels were built loose to allow for thermal expansion, and looked --- from
below, at that altitude, in sunlight --- like nothing anybody had a category for. The U-2 flew
above 70,\!000 feet from 1955, at a time when commercial aviation topped out around 20,\!000 and
airline crews were reporting silver objects far above them. The CIA's own declassified history of
the U-2 programme states plainly that the aircraft accounted for a substantial fraction of UFO
reports in the late 1950s, and that Blue Book investigators were quietly given enough
information to dismiss those reports without explaining why.

That is the single most important fact in this chapter. There is a documented, admitted case of
the US government using the UFO phenomenon as cover for a classified programme. It is not a
theory. And it cuts both ways: it vindicates the conspiracy-minded about the \emph{behaviour},
and it removes the anomaly from the \emph{sightings}.

\begin{funfact}
The F-117 Nighthawk flew in 1981 and was not acknowledged until 1988. For seven years there were
faceted, matte-black, impossibly angular aircraft operating at night out of Tonopah, and people
saw them. The Nevada triangle-UFO reports of the 1980s look considerably less mysterious once you
know the flight schedules --- and the B-2, first seen publicly in 1988, is a large silent-seeming
flying wing that from directly below is a black triangle. Chapter~\ref{ch:ufocsi}'s Phoenix Lights deserve to be
read with that in mind.
\end{funfact}

% ---------------------------------------------------------------
\section{Theory of Flight}\label{sec:theoryflight}

\hi{The question this section has to answer is not whether UFOs can fly.} That one is settled, and
it was settled in front of a congressional committee. On 9 September 2025 the House task force on
unidentified anomalous phenomena was shown gun-camera footage, released by Representative Eric
Burlison, of an MQ-9 Reaper tracking a small luminous object off the coast of Yemen on 30 October
2024 while a second Reaper put an AGM-114 Hellfire into it. The object was still there afterwards
and left the frame under its own power.

I want to be careful with that tape, because the version circulating online has the orb
\emph{shooting the missile down}, and it does not. What it shows is a strike --- or a very near
miss --- followed by an object carrying on as though nothing had happened. Neither the Department
of Defense nor AARO has confirmed the provenance, the resolution is poor enough that a fusing
failure or a parallax miss stays on the table, and I would not build a physics on fourteen seconds
of infrared. But a Hellfire arrives at better than 1,\!000 miles an hour carrying roughly twenty
pounds of high explosive, and if the footage is what it appears to be, then the interesting
question is not how the thing flew away afterwards. \hi{It is what there was to hit.}

Hold that, because everything below assumes an answer to it.

You cannot judge whether an object's flight is impossible without knowing what makes flight
possible. This section is the physics, briefly, because it is the tool you need for Chapter~\ref{ch:ufocsi}'s
video analysis and for every ``it performed a 90-degree turn at 7,\!000 miles an hour'' claim.

Four forces act on an aircraft: lift, weight, thrust, drag. Lift is generated by the wing turning
the airflow downward; by Newton's third law, the reaction is upward. The popular ``longer path
over the top'' explanation is wrong, and demonstrably so --- it predicts the wrong magnitude and
fails to explain inverted flight. Circulation theory, via Kutta--Joukowski, gets it right.

The number that matters is the one nobody mentions: \textbf{$g$-loading}. A turn is an
acceleration, and the acceleration required scales as velocity squared over radius. \hi{This is why
the reported manoeuvres are the real problem, not the speeds.}

Run it. An object at 3,\!000 miles per hour --- roughly 1,\!340 metres per second --- making a
turn with a radius of 100 metres experiences an acceleration of $v^2/r$, which is about
18,\!000,\!000 m/s$^2$, or around 1.8 million $g$. Human tolerance is roughly 9 $g$ sustained
with a suit and training; the highest survived brief peak on record is around 46 $g$ (John Stapp,
who did it to himself on a rocket sled and detached his retinas). Structural limits for a fighter
airframe are in the region of 9 to 12 $g$. Even an unmanned airframe of any known material fails
three or four orders of magnitude below what these reports require.

The consequence is important and it is not the one either camp usually draws. If an object really
performed those manoeuvres, the occupants and the structure would need to be shielded from
inertia itself --- which means the interesting claim is not propulsion but inertial mass
manipulation, and no known physics offers a route to it. \emph{Or} the reported values are wrong,
which happens the moment you misjudge range: angular motion across your field of view is
unremarkable for a nearby small object and impossible for a distant large one, and observers
cannot judge range on an unfamiliar object against an empty sky. That is not a slur on witnesses.
\hi{It is a documented limitation of monocular distance estimation beyond a few hundred metres.}

Two more numbers worth carrying. \textbf{Sonic boom}: anything exceeding roughly 767 mph at sea
level produces a shock wave that is heard for miles. \hi{Silent supersonic flight in atmosphere is
not a stealth feature; it is a violation of gas dynamics.} \textbf{Heating}: at Mach 3 the SR-71's
leading edges reached over 300$^\circ$C, which is why it was built of titanium. An object doing
Mach 20 in the lower atmosphere would be incandescent and would ablate.

\begin{funfact}
The hardest number for the exotic-craft hypothesis is not any of these. It is the absence of
turbulence. Every one of those manoeuvres displaces an enormous mass of air, and displaced air is
visible on radar as a disturbance, audible as noise, and detectable by aircraft flying behind ---
wake turbulence is why air traffic control enforces separation minima. Objects reported to shift
from hover to hypersonic leave no wake, no boom, no thermal plume and no downwash on the water
beneath them. In physics, if you move something you have to move what it was in.
\end{funfact}

\subsection{The Solidity Assumption}\label{sec:solidity}

Every number in this section rests on a premise nobody in this field states out loud, including
me, up to now. \hi{We have been assuming the object is solid.} $g$-loading is a statement about a
rigid body with mass and a structure to fail. Wake turbulence assumes a hull displacing air.
Ablation assumes a leading edge made of something. Take solidity away and the whole apparatus
above --- my apparatus, the one I have just spent four pages sharpening --- stops applying.

And the archive has been telling us to take it away for seventy years. Look at the catalogue in
Section~\ref{sec:classificationcommonufos}. The outlines that recur most stubbornly are not
outlines at all: a disc that becomes a sphere, a cylinder that divides into three and rejoins, a
formless glow with no fixed edge between one glance and the next. My taxonomy has a card for the
residual case because the reports demanded one. \hi{A rigid airframe has exactly one shape.}
Whatever else a shape change is, it is not a manoeuvre, and it cannot be flown.

So take the ordinary explanations first, as ever, and they get most of it. A disc seen edge-on is
a cigar; a tumbling object presents a sequence of silhouettes; a luminous blob at the limit of
resolution changes shape because your eye is inventing an edge that the photons never supplied.
Much of the transformation literature is geometry and optics rather than physics, and I would put
the majority of the shape-shifting corpus there without apology. What survives that filter
points somewhere uncomfortable: at plasmas, at ionised air, at luminous phenomena with no surface
to touch --- things that would have no $g$-limit to violate because they have no structure to
break, and would leave no wake because there is nothing much there to displace.

Which leaves one question, and it is answerable, because the corpus contains precisely one
relevant measurement. In eighty years of this subject, outside a laboratory and outside anybody's
classified programme, \emph{one} man has put his hand on one of these objects and lived to file a
report. Stefan Michalak, Falcon Lake, Manitoba, 20 May 1967: a hull hot enough to melt the
fingertips of his glove, a panel of vents, and a jet of gas that set his shirt on fire and left
him ill for over a year. That case is in Chapter~\ref{ch:nuclear}, Section~\ref{sec:falconlake},
because it is the case the nuclear-propulsion argument leans on hardest; but it belongs here just
as much, because it is the only tactile datum we have.

\hi{And it says solid.} Hot metal, machined openings, mechanical exhaust --- an object with a
surface, obeying the physics in this section, made by somebody. One data point, uncontrolled,
sixty years old, from a witness with a sceptical case against him that I have set out in full
elsewhere and do not withdraw. That is what the entire physical basis for the solidity assumption
rests on: $n = 1$, and it nearly killed him. Anybody who tells you UFOs are solid craft is
extrapolating from Michalak's glove, and anybody who tells you they are immaterial has to explain
the grid burned across his stomach. \hisoft{I hold both, uncomfortably, and I would rather report
the discomfort than resolve it for you.}

% ---------------------------------------------------------------
\section{Foo Fighters}\label{sec:foofighters}

During the Second World War, Allied aircrews --- particularly night fighter crews of the US 415th
Night Fighter Squadron over the Rhine valley in late 1944 --- reported small glowing balls of
light, usually orange or red, that paced their aircraft, held formation, appeared to manoeuvre
intelligently, and did not attack. The crews named them foo fighters, from the ``Smokey Stover''
comic strip catchphrase ``where there's foo, there's fire''. Japanese and German crews reported
comparable things in their own theatres, which is the detail that makes the case robust: these
were not one air force's rumour.

The reports were taken seriously precisely because the obvious hypothesis was enemy technology,
and each side suspected the other. Neither had it. That mutual suspicion, incidentally, is a
general lesson about secret weapons: the people who would know are usually the ones asking.

The candidate explanations are several, and unusually good.

\textbf{St.~Elmo's fire} --- corona discharge from a charged airframe --- is well known to
aircrew, appears at propeller tips, wingtips and windscreen edges, glows blue or orange, and
moves. It is common in the electrical conditions of the northern European winter.

\textbf{Ball lightning} is real, now photographed and reproduced under some conditions, and
remains partially unexplained --- luminous spheres persisting for seconds, associated with
thunderstorms, occasionally reported entering aircraft.

\textbf{Ice crystal and lens artefacts}, reflections from a windscreen, and internal reflections
in gunsight optics account for a category of paced lights, and the \emph{pacing} is the tell: an
object that maintains perfect station with your aircraft regardless of your manoeuvre is very
likely attached to your aircraft, either physically or optically.

\textbf{Autokinesis} makes a single point light in a dark field appear
to drift and dart. Night fighter crews were staring into blackness for hours, which is the exact
configuration in which this occurs.

\textbf{And some were real hardware.} German flak used pyrotechnic tracking rounds and there were
experimental programmes; V-1s and V-2s were airborne; and the Me 163 rocket interceptor and Me 262
left unfamiliar signatures at night.

What I find valuable about the foo fighters is not the identification. It is the timing. These
reports predate Kenneth Arnold in 1947, and they are noticeably \emph{different} in content:
lights that pace, not discs that fly. Once the flying saucer template existed, the reports
changed shape. That is a cultural fingerprint on raw perception, visible in the archive, and it is
Section~\ref{sec:mufon}'s argument with a date stamp.

% ---------------------------------------------------------------
\section{The Credibility of Pilots}\label{sec:pilots}

The most common appeal in this field is to the authority of the pilot: trained observer,
thousands of hours, no reason to lie, therefore his report should be weighted heavily. The appeal
is reasonable, and it is also substantially wrong, and explaining why requires distinguishing two
things that get conflated.

Pilots are highly reliable about \emph{what they did, what their instruments read, and what
procedures they followed}. That is what their training is for and their record is excellent.

Pilots are \emph{not} reliable at judging the size, range, or speed of an unfamiliar object
against an empty sky, and there is no training that makes anybody reliable at it, because the
necessary visual cues are absent. This is not my opinion; it is the consistent finding of the
aviation human-factors literature, and it is why the industry built an entire safety culture
around not trusting the eyeball.

The specific failure modes:

\textbf{The empty-field problem.} With no reference objects, the eye has no basis for range, and
the brain therefore assigns one from assumed size --- circularly. A weather balloon at 20 miles
and a metallic disc at 2 miles produce identical retinal images.

\textbf{Relative motion and parallax.} A distant object appears to move with your aircraft; a
nearby one appears to shoot past. This is why Venus has been chased by fighters --- repeatedly and
embarrassingly --- and why a celestial body appears to pace, then dart away when the pilot turns.

\textbf{Vection and false horizons.} At night, over water, or in cloud, the vestibular system
misreports attitude confidently, which is what killed John Kennedy Jr. Pilots are trained
explicitly to disbelieve their own sensations and fly the instruments, which is an institutional
admission of exactly the point I am making.

\textbf{Canopy and HUD artefacts.} Reflections, scratches, coatings, and symbology produce
apparent objects with apparent motion.

\textbf{Workload and attention.} A crew running an intercept has extremely high cognitive load,
and load degrades peripheral perception and memory encoding.

\figwrap[r]{0.38}{Kenneth Arnold with his own drawing of the objects seen near Mount Rainier --- crescents, not saucers.}{https://commons.wikimedia.org/wiki/File:Kenneth_Arnold_UFO_drawing.jpg}{fig:arnold}
The key illustration is Kenneth Arnold, the founding case. On 24 June 1947 he was a private pilot
with a legitimate reputation, flying near Mount Rainier (see Figure~\ref{fig:arnold}), and he reported nine objects moving in a
way he described as ``like a saucer if you skipped it across water''. His \emph{motion} simile
became ``flying saucer'' in the press --- he did not describe the objects as saucer-shaped, and
he complained about the misquotation for the rest of his life. His speed estimate of over 1,\!200
mph was derived by assuming a size he had no way to know. The most likely candidates are a
formation of pelicans catching sunlight, or lenticular cloud fragments. And the entire modern era
of this subject grew from a misreported simile.

\begin{funfact}
The airlines eventually solved the pilot-credibility problem in a way ufology never has, and it
is the best argument in this book for how to handle testimony. After a series of accidents caused
by confident, senior, experienced men being wrong, aviation adopted Crew Resource Management,
mandatory cross-checking, instrument primacy, and the non-punitive Aviation Safety Reporting
System run by NASA specifically so that reports could be filed without career risk. Commercial
aviation became the safest form of travel ever devised by systematically distrusting individual
perception. Ufology's response to the same problem was to add ``he was a very experienced pilot''
to the beginning of the sentence.
\end{funfact}

\subsection{The best-credentialled witness in ufology: John Lear}\label{sec:johnlear}

If the argument above is right, then the field's strongest possible pilot-witness should also be
its clearest demonstration of why credentials do not settle anything. That witness exists, and he
is Bill Lear's son.

John Olsen Lear (1942--2022) held an Airline Transport Pilot certificate and every category and
class rating the Federal Aviation Administration issues, plus type ratings in a couple of dozen
aircraft. He claimed something over 19,\!000 hours, flew for more than twenty operators across the
cargo and charter world between 1967 and 1983 --- including contract work that he later described,
and that others have described, as flying for the Central Intelligence Agency in Southeast Asia,
Africa and Central America --- and in 1966, at twenty-four, he set a round-the-world speed record
for light jets in one of his father's aircraft: roughly 22,\!000 miles in 50 hours and 39 minutes
of flying time. He was, on paper, exactly the man the appeal-to-the-pilot argument is built for.
He was also disinherited, and he spent the second half of his life in Nevada rather than in the
family business.

\figwrap[r]{0.38}{The Groom Lake Road boundary: warning signs, a camera mast and a white truck on the ridge. This is the frontier where Lear's story and Lazar's converge.}{https://commons.wikimedia.org/wiki/File:Warnhinweis_und_Sicherheitspersonal_an_der_Groom_Lake_Road_07.2008.jpg}{fig:groomroad}
What he did with that authority is the point. In December 1987 he posted a statement to the
ParaNet bulletin board asserting that the government had recovered crashed craft and living
occupants, that a treaty had been signed, and that there was a joint human--alien facility under
Dulce, New Mexico, at which things were being done to abducted people that he described in the
most lurid terms available. There was no documentation behind any of it that survives inspection;
the material traces back to the MJ-12 papers and to the Krill correspondence discussed in
Chapter~\ref{ch:underground}, which is to say to a chain of anonymous documents rather than to anything
Lear saw from a cockpit. He became MUFON's Nevada state director, spoke at the July 1989
symposium, and lost that position in the same period, because a state directorship is a
volunteer job in an organisation that has to worry about its own credibility.

And then, in March 1989, he introduced a Las Vegas television reporter named George Knapp to a
young man who said he had worked on flying discs south of Groom Lake. Everything in
Section~\ref{sec:s4lazar} --- Element 115, S4, the sport model, the whole modern architecture of
Area 51 --- reaches the public through that introduction. Lear is the load-bearing member. Take
him out and Bob Lazar is a man with a story and no microphone.

The two of them also went to look. On the night of 5--6 April 1989 a group that included Lazar,
Lazar's wife, her sister, Lear and a fifth man named John Huff drove out to a ridge off Groom Lake
Road to watch and film what Lazar said were test flights on a predictable Wednesday schedule (see
Figure~\ref{fig:groomroad}). They were stopped by security, held, and turned over to the Lincoln
County Sheriff's Office. I flag the provenance honestly: the incident is documented in
researchers' accounts and in Lear's and Lazar's own later retellings rather than in any primary
document I have been able to obtain, and the various versions disagree on who was detained, what
was confiscated and what was charged. Treat the arrest as reported, not as established.

So hold the two halves together. Nineteen thousand hours, every rating on the certificate, a
world record and a famous surname on one side; MJ-12, Dulce, and a treaty with the Greys on the
other. The credential is real and it purchases nothing, because the claims were never about what
Lear observed --- they were about documents he was handed and believed. That is the whole thesis
of this section in one biography, and it is why the sober response to ``he was a very experienced
pilot'' is ``experienced at what, exactly?'' Lear later took his views to the radio, where
Section~\ref{sec:artbell} picks the thread up again.

% ---------------------------------------------------------------
\section{Aliens At The Airport}\label{sec:airport}

Airports ought to be the best UFO laboratories on the planet, and it is worth understanding why
they so consistently disappoint.

A major international airport is a volume of monitored airspace with primary and secondary radar,
multilateration, ADS-B, surface movement radar, thousands of pairs of trained eyes, dozens of
fixed cameras, and a legal obligation to write things down. If unknown craft were routinely
operating in our atmosphere, the terminal areas of Heathrow, O'Hare, Haneda and Frankfurt are
where we would have caught them decades ago, in calibrated detail, with the tapes preserved. What
we have instead is a small number of famous airport incidents whose evidentiary record is
startlingly thin --- and the reasons for the thinness are themselves instructive.

\hi{Take Chicago O'Hare, 7 November 2006, the most cited airport case in the modern literature.}
Around 4:15 in the afternoon, a United Airlines ramp employee at Gate C17 reported a dark grey
disc hovering perhaps 1,\!900 feet above the terminal, motionless, silent. Word went round the ramp
frequency. Over the following minutes somewhere between a dozen and twenty United employees ---
ramp staff, a supervisor, at least one pilot --- reported seeing it, and several said it departed
vertically at speed, punching a clean hole through the overcast that stayed visible for a short
while afterwards. That last detail is the case's whole reputation: a hole in a cloud deck is a
physical effect, and physical effects can in principle be tested.

Here is what the record actually holds. United initially denied any knowledge of the event; the
\emph{Chicago Tribune} obtained the internal reports and the FAA's initial position via a Freedom
of Information Act request, at which point both organisations confirmed employees had filed
accounts. The FAA's formal conclusion was a weather phenomenon --- a hole-punch cloud --- and it
conducted no investigation beyond that. No O'Hare radar recorded a target. Nobody, in an airport
handling roughly a quarter of a million passengers a day, in a terminal ringed with cameras and
full of people carrying cameraphones in 2006, produced a single photograph. Multiple witnesses
described the hole in the cloud; not one documented it.

\hisoft{The cloud hole is the part sceptics should engage with honestly rather than wave away.}
Hole-punch clouds and fallstreak holes are real, well-photographed meteorological features: a
supercooled altocumulus layer glaciates locally, the ice crystals fall out, and a startlingly
crisp circular gap opens. They are frequently triggered by aircraft climbing through the layer,
which is not a small coincidence at an airport. The honest reading of O'Hare is that a genuine
atmospheric feature was observed by people with no reason to lie, in a working environment where
the interpretive frame ``something took off through the cloud'' arrives free of charge --- and
that because nobody preserved a single instrument record or image, there is now no way to test
any of it. The witnesses were credible. The evidence was never collected. That combination is the
central pathology of this entire subject, and it is not solved by the witnesses being credible.

The other much-repeated case runs the opposite way. On 7 July 2010, Hangzhou Xiaoshan
International Airport in China closed for around four hours and diverted eighteen aircraft after
an unidentified object was reported in the approach airspace. Here the institutional response was
everything ufology says it wants: the airspace was actually shut, the event was reported by state
media, and a photograph exists. The closure is the fact of the case --- and it establishes
precisely nothing about the object, because closing airspace over an unidentified radar or visual
contact is what a competent authority is supposed to do regardless of what the contact turns out
to be. Chinese authorities subsequently indicated the object was related to military activity and
declined to elaborate; the widely circulated photograph is a long-exposure of a light streak
entirely consistent with an aircraft or a rocket stage catching sunlight. Hangzhou is a case about
airspace management, not about visitation.

\begin{funfact}
Gatwick, December 2018, is the control experiment nobody in ufology wants to run. Britain's
second-busiest airport shut for roughly 33 hours, cancelling about 1,\!000 flights and stranding
140,\!000 passengers, on the strength of around 170 reported sightings of drones over the
airfield --- from police officers, airport staff and pilots. Two people were arrested and
released without charge. After the largest investigation of its kind, Sussex Police could not
confirm that any of the 170 sightings involved an actual drone, and acknowledged that some may
have been police drones or misidentified aircraft lights. Same setting, same class of trained
observers, far better documentation, an object type nobody disputes exists and can be bought in a
shop --- and the sighting record still could not be resolved. If 170 witnesses cannot reliably
establish a drone, we should be extremely modest about what twenty witnesses establish regarding
a spacecraft.
\end{funfact}

There is a broader pattern here, and it cuts against the visitation hypothesis rather than for it.
Airport cases cluster around the ramp and the tower, not the runway; they produce recollections
rather than recordings; and they almost never appear in the radar data that would settle them in
an afternoon. Meanwhile the genuine mystery traffic at airports --- drones, laser strikes,
balloons, sky lanterns, misidentified Venus, and the occasional untracked light aircraft --- is
logged in enormous, boring quantity, and is investigated, and is usually resolved. The signal in
that asymmetry is hard to escape. The environments with the best instrumentation on Earth generate
the fewest good UFO cases. That is exactly the distribution you would expect if the phenomenon
were principally perceptual, and exactly the opposite of what you would expect if craft were
visiting.

% ---------------------------------------------------------------
\section{UFOs \& UAP}\label{sec:uap}

Somewhere around 2017 the vocabulary changed, and the change matters more than it looks.
\term{UFO} --- unidentified flying object, coined by Captain Edward Ruppelt of Project Blue Book in
1952 precisely because ``flying saucer'' had become unusable --- had by the 1990s acquired the same
problem. Say ``UFO'' in a briefing room and everyone hears \emph{alien}, which is exactly the
inference the term was invented to avoid. So the U.S.\ Navy adopted \term{UAP}: originally
Unidentified Aerial Phenomena, broadened in 2022 to Unidentified \emph{Anomalous} Phenomena when
the Pentagon's remit was extended to objects in space and underwater as well as in the air. The
official body is now AARO, the All-domain Anomaly Resolution Office, which succeeded the UAP Task
Force, which succeeded AATIP, the Advanced Aerospace Threat Identification Program.

I am mildly cynical about the rebrand and I will say why: it is a laundering operation for
respectability, and it worked. Identical claims that were career-ending in 1995 became briefable to
the Senate Intelligence Committee in 2019 without any new evidence being produced --- only a new
noun. That is worth noticing as a lesson about institutions rather than about physics. But I also
think the term is genuinely better, because it drops the two assumptions smuggled inside
``flying object'': that the thing is flying, and that it is an object. A sensor artefact is neither.

\subsection{The Nimitz Incident}

On 14 November 2004, aircraft from Carrier Strike Group Eleven, operating with USS \emph{Nimitz}
about a hundred miles southwest of San Diego, were vectored onto an unknown contact. The cruiser USS
\emph{Princeton}, running the then-new SPY-1B radar upgrade, had been tracking objects intermittently
for several days: appearing at around 80,000 feet, descending to roughly 20,000 feet, and holding.

Commander David Fravor, then commanding officer of the Black Aces --- Strike Fighter Squadron 41 ---
was airborne on a training sortie with his wingman when \emph{Princeton} tasked him with a real-world
intercept. What he reported seeing, and has reported consistently for two decades in terms that have
not inflated, was a disturbance in the sea surface roughly the size of a Boeing 737, and above it a
white object about forty feet long, shaped like an elongated Tic Tac, with no wings, no rotors, no
exhaust plume and no visible means of propulsion, moving erratically over the water. As he descended
to engage it, he says it rose to meet him, then accelerated away in a manner he could not match.
\emph{Princeton} then reported reacquiring a contact at the CAP point sixty miles away --- the
squadron's pre-briefed rendezvous, which is the detail that makes the case unsettling rather than
merely odd.

\figwrap[l]{0.44}{A frame from the 2004 ``Tic Tac'' ATFLIR capture made during the
\emph{Nimitz} encounter, formally released by the Department of Defense in April 2020.}{https://commons.wikimedia.org/wiki/File:FLIR1_Official_UAP_Footage_from_the_USG_for_Public_Release.webm}{fig:flir1}
A later flight that day captured the FLIR footage subsequently released as ``FLIR1'' or the ``Tic
Tac'' video. Fravor was not the pilot who filmed it. In 2017 the video appeared in \emph{The New York
Times} alongside the disclosure of AATIP, and in April 2020 the Department of Defense formally
released it and two other clips, confirming they were genuine Navy recordings and that the objects
remained unidentified.

\subsection{Why This Case Is Different, and Where It Is Weak}

Different, first. Fravor is about as good a witness as this subject will ever get: a Navy commander
with roughly 3,000 flight hours, TOPGUN-qualified, who reported through the chain of command at the
time, gained nothing and risked a good deal, and whose account has not grown in the telling. There
were multiple independent sensor systems --- shipboard radar, airborne radar, infrared --- and
multiple aircrew across two flights. Section~\ref{sec:pilots} argued that pilot credibility is
oversold in general; this is the case that survives that argument best, because the claim does not
rest on his eyes alone.

Now the weaknesses, because they are real and the enthusiast literature omits them. The
\emph{Princeton}'s radar had just been upgraded and was reportedly generating spurious tracks, which
is exactly the profile of a new-software artefact. The video itself has been analysed in detail by
Mick West and others, and the ``impossible'' motion in it is largely explicable: the apparent
rotation is the gimbal of the ATFLIR pod rolling as it tracks, and the apparent high-speed departure
at the end coincides with the operator changing zoom level, whereupon a distant object at the edge of
tracking lock simply leaves frame. Those are not hand-waves; they are demonstrated with the
instrument's own documented behaviour. And critically, the compelling parts of Fravor's account ---
the water disturbance, the ascent, the acceleration --- are \emph{not} in the video at all. The
footage everyone argues about is not the encounter everyone finds impressive.

So my position: the Nimitz case is the strongest single item in the modern corpus and it still gives
us radar returns from an instrument with a known fault, an infrared clip with a plausible mundane
reading, and unverifiable testimony --- excellent testimony, but testimony. AARO's 2024 historical
review found no evidence of extraterrestrial technology in any case it examined, including this one,
while leaving a residue of genuinely unresolved reports. That residue is where this subject actually
lives, and Section~\ref{sec:deconstructingpeerreview} is why I will not fill it in with a guess.

\begin{funfact}
The 2021 Office of the Director of National Intelligence preliminary assessment examined 144
incidents from 2004--2021 and managed to explain exactly one --- a deflating balloon. That statistic
is quoted constantly as ``143 unexplained''. What the report actually said is that most cases lacked
sufficient data to reach \emph{any} conclusion. ``We could not determine what it was'' and ``it was
something extraordinary'' are different sentences, and the entire modern disclosure debate runs on
blurring them.
\end{funfact}

\begin{srcnotes}
\item Cooper, H., Blumenthal, R.\ \& Kean, L.\ (16 December 2017). Glowing Auras and Black Money.
\emph{The New York Times}.
\item Office of the Director of National Intelligence (25 June 2021). Preliminary Assessment:
Unidentified Aerial Phenomena.
\item All-domain Anomaly Resolution Office (2024). Report on the Historical Record of U.S.\
Government Involvement with UAP, Volume I.
\end{srcnotes}

% ---------------------------------------------------------------
\section{The Buga Sphere}\label{sec:bugasphere}

Everything in Section~\ref{sec:uap} is at least twenty years old, which lets both sides argue from
memory. The Buga sphere is the opposite: it is the most recent case in this book, it is still being
worked on week by week as I write, and for once the disputed object is not a smear of infrared pixels
but a physical thing sitting on a bench that somebody can weigh. That makes it the best available
test of whether this field can actually run a physical investigation --- and, so far, of how badly it
runs one.

On 2 March 2025, according to the men who brought it forward, a metallic sphere was seen manoeuvring
over Buga, a town in the Valle del Cauca department of western Colombia. The account is that it
struck power lines near the Guadalajara river and came down, and that three local men recovered it.
What has since been shown to cameras and to a handful of researchers is a dull metallic ball roughly
the size of a basketball, weighing on the order of two kilograms --- about four and a half pounds ---
with fine incised markings on its surface that the promoters describe as engraved symbols or glyphs.
Nothing about that description is by itself remarkable. Everything remarkable about the case is in the
claims attached to it.

\subsection{The claims made for the object}

The physical testing was arranged by the group promoting the sphere rather than by any independent
body. Jos\'e Luis Vel\'asquez, leading a Mexican team, has said the object has no welds, no seams and
no joints of any kind, and that X-ray imaging shows an internal structure of three concentric layers
enclosing a set of small free-floating spheres. How many small spheres is a question the case has
never answered consistently: nine in some accounts, eighteen in others, from the same investigative
group, in the same year. A trivial, directly countable fact --- the sort of thing a single published
radiograph with a scale bar would settle permanently --- has instead drifted, and no radiograph has
been released for anyone else to count.

Julia Mossbridge, of the University of San Diego, was among the small number of credentialled
researchers given physical access. Her published reaction is worth quoting for its restraint: the
object, she said, ``looks so human made,'' and her suggestion was that it might be somebody's ``cool
art project.'' Her recommendation was the correct one and it has not been followed --- hand the sphere
to a group set up to do exactly this work, such as the Galileo Project or the Scientific Coalition for
UAP Studies, and let them cut it, image it, and publish. Instead the object has stayed with the people
whose interests depend on what it turns out to be.

There is a further problem with the discovery footage, which is the material most of the world actually
saw. Skeptical investigators, chiefly the contributors at Metabunk, traced the videos to promotional
material connected with a Colombian retailer of metal-detecting equipment --- a company whose own logo
is visible on the backdrop of the ``laboratory'' in which the sphere is examined, and whose product
line advertises detection technology combined with dowsing. The flight sequences are consistent with a
light ball flown on a line beneath a drone, or with a weighted balloon, and in at least one clip the
object leaves frame and is reacquired travelling in a different direction. None of that is proof of
fraud in the strict sense, and I want to be careful here: an object can be genuinely anomalous and
still be promoted by people making dreadful videos. But it does mean the provenance of the footage is
commercial, and provenance is the whole ballgame in physical-evidence cases.

\subsection{The radiocarbon claim}

In September 2025 the case acquired the sort of number that travels well. Steven Greer announced that
the University of Georgia's Center for Applied Isotope Studies had accelerator-mass-spectrometry dated
resin taken from thirty-one microscopic holes in the sphere to roughly 12,560 years before present ---
which, as every promoter of the story immediately noticed, drops the object into the Younger Dryas.
An artefact from the exact geological window that already carries a heavy load of lost-civilisation
speculation is a marketing gift, and it was treated as one.

Take the report at face value for a moment and it still does not say what it is being used to say.
The document names its sample material inconsistently, describing a foraminifer sample in one place
and a resin in another; foraminifera are single-celled marine organisms whose carbonate shells are the
standard workhorse of ocean-sediment dating, and the pretreatment described --- a hydrochloric acid
wash to strip carbonates --- would destroy the very shells the header claims were dated. Nobody
outside the promoters' circle has been told where on the object the sample was taken. And even in the
most generous reading, dating a speck of carbonate-bearing sediment lodged in a hole tells you the age
of the sediment, not the age of the ball it is stuck to. Ancient forams adhering to a modern surface
make the forams old. They do not make the sphere old, any more than a Roman coin found in a coat
pocket ages the coat.

This is the pattern I have complained about for nineteen chapters, so I will name it plainly rather
than pretend the Buga case is a special disappointment. A number with a laboratory's letterhead on it
was released by an interested party, through a personal assistant rather than through a journal, with
no sampling record, no stratigraphic context and no chain of custody --- and within days it had become,
in popular retelling, ``scientists date UFO to 12,500 years ago.'' The laboratory did the small thing
it was asked to do. Everything that would make the result mean anything happened, or failed to happen,
outside the laboratory.

\subsection{What would settle it}

I am going to be unusually concrete, because this is one of the rare cases where the decisive
experiment is cheap and nobody is doing it. Give the sphere to a laboratory with no stake in the
outcome. Publish a computed-tomography scan, so the layer count and the sphere count stop being a
matter of testimony. Run a full elemental and isotopic analysis of the shell metal itself --- not of
debris in its crevices --- because isotope ratios in industrially refined metal are a fingerprint of
terrestrial ore processing and would end the argument in an afternoon. Section a small coupon and put
it under an electron microscope: manufacturing leaves grain structure, and grain structure does not
lie about whether something was cast, spun, sintered or additively printed. Then publish the lot,
including the sampling map, and let other people repeat it.

Until that happens my assessment is simple, and it is not the exciting one. The most probable
explanation of the Buga sphere is that it is a terrestrial artefact --- quite possibly an art object or
a deliberately made curiosity --- attached to a commercial promotion, and that the anomalous readings
reported for it are artefacts of testing arranged by people who needed a particular answer. I hold that
position loosely, because unlike almost every other case in this chapter it is falsifiable this year by
anyone granted an hour with the object and a CT scanner. That is a genuinely unusual situation in
ufology, and the fact that the opportunity is being wasted is itself the most informative datum the
case has produced.

\begin{funfact}
The Buga sphere is one of the very few UFO cases in history where the object in dispute is available,
in one piece, on a table, and the answer could be had in a single afternoon of destructive testing.
Six months after it was recovered, the thing that had not been done was the testing. Compare the
Roswell debris, the Ubatuba magnesium, the Falcon Lake soil samples: in seventy-eight years the field
has repeatedly asked for physical evidence, and on the rare occasions it gets some, it takes it on
tour instead.
\end{funfact}

\begin{srcnotes}
\item Center for Applied Isotope Studies, University of Georgia (19 September 2025). Radiocarbon
dating report issued to Sirius Technology Advanced Research LLC; circulated by Steven Greer.
\item Metabunk (2025). ``Recently viral Buga, Colombia, `alien' metal balls'' and ``Buga Sphere
Carbon Dating'' investigation threads.
\item Mossbridge, J.\ (2025). Public statements following physical examination of the Buga sphere.
\item No photograph of the sphere with a verifiable free licence and confirmed provenance was
available at the time of writing, so this section is deliberately unillustrated rather than
illustrated with a still lifted from the promotional footage discussed above.
\end{srcnotes}
